% Copyright 2023 Sycuricon Group
% Author: Phantom (phantom@zju.edu.cn)

\section*{Problem 3}

Consider the following two 9-bit floating-point representations based on the IEEE floating point format.

\begin{itemize}
    \item Format A: There is 1 sign bit. There are k = 5 exponent bits. The exponent bias is 15. There are n = 3 fraction bits.
    \item Format B: There is 1 sign bit. There are k = 3 exponent bits. The exponent bias is 3. There are n = 5 fraction bits.
\end{itemize}

\noindent In the following table, you are given some bit patterns in format A, and your task is to convert them to the value in format B.
In addition, give the values of numbers given by the format A and format B bit patterns.
Give these as whole numbers (e.g., 17) or as fractions (e.g., 17/64).
If the value represented by bits in format A cannot be expressed by format B, please fill \textbf{NULL} in the \textbf{Value} column and explain the reason briefly in \textbf{Bits} column.

\begin{table}[h]
    \center
    \begin{tabular}{c|>{\centering}p{3cm}|p{3cm}|p{3cm}|p{3cm}}
    \hline
        & \multicolumn{2}{c|}{Format A} & \multicolumn{2}{c}{Format B} \\ \hline
    No. & \multicolumn{1}{c|}{Bits} & \multicolumn{1}{c|}{Value} & \multicolumn{1}{c|}{Bits} & \multicolumn{1}{c}{Value} \\ \hline
    1   & 0 10010 011 &       &      &       \\
    2   & 1 00011 010 &       &      &       \\
    3   & 0 00011 010 &       &      &       \\
    4   & 1 11000 000 &       &      &       \\
    5   & 0 10011 100 &       &      &       \\ \hline
    \end{tabular}
\end{table}

\subsection*{Answer:}
